% chapters/04_ergebnisse.tex
% Viertes Kapitel: Ergebnisse
% Stellt die empirischen Befunde der Untersuchung dar.

\chapter{Ergebnisse}
\label{chap:ergebnisse}

Dieses Kapitel praesentiert die Ergebnisse der in Kapitel~\ref{chap:methodik}
beschriebenen Untersuchung. Die Darstellung folgt der Reihenfolge der
Forschungsfragen.

\section{Ergebnisse zu Forschungsfrage 1}
\label{sec:ergebnis_ff1}

Hier werden die Befunde zur ersten Forschungsfrage beschrieben. Zentrale
Ergebnisse koennen durch Tabellen oder Abbildungen veranschaulicht werden.

% Abbildung einbinden (Datei im figures/-Ordner ablegen):
% \begin{figure}[h]
%   \centering
%   \includegraphics[width=0.75\textwidth]{figures/ergebnis-ff1.png}
%   \caption{Ergebnis der Auswertung zu Forschungsfrage 1}
%   \label{fig:ergebnis-ff1}
% \end{figure}

Tabelle~\ref{tab:ergebnisse_ff1} fasst die quantitativen Kennzahlen zusammen.

\begin{table}[h]
\centering
\caption{Ergebnisse zu Forschungsfrage 1}
\label{tab:ergebnisse_ff1}
\begin{tabular}{lcc}
\toprule
\textbf{Kennzahl} & \textbf{Wert} & \textbf{Einheit} \\
\midrule
Kennzahl A        & 42{,}3        & \%               \\
Kennzahl B        & 17            & Anzahl           \\
Kennzahl C        & 3{,}8         & Punkte           \\
\bottomrule
\end{tabular}
\end{table}

\section{Ergebnisse zu Forschungsfrage 2}
\label{sec:ergebnis_ff2}

Hier folgt die Darstellung der Befunde zur zweiten Forschungsfrage.

\section{Diskussion der Ergebnisse}
\label{sec:diskussion}

Die Ergebnisse werden in diesem Abschnitt im Licht der bestehenden Literatur
eingeordnet. Dabei wird untersucht, inwieweit die Befunde mit den Aussagen
von \textcite{beispiel_artikel_2019} uebereinstimmen oder davon abweichen.

Moegliche Erklaerungen fuer abweichende Befunde sind:
\begin{enumerate}
  \item Erklaerung 1 \ldots
  \item Erklaerung 2 \ldots
\end{enumerate}

\section{Limitationen}
\label{sec:limitationen}

Jede empirische Untersuchung unterliegt methodischen Grenzen. Die wesentlichen
Limitationen dieser Arbeit sind:

\begin{itemize}
  \item \textbf{Stichprobengröße:} Die Stichprobe umfasst nur \ldots,
        was die Verallgemeinerbarkeit der Ergebnisse einschraenkt.
  \item \textbf{Erhebungszeitraum:} Die Daten wurden im Zeitraum \ldots
        erhoben und koennen saisonale Schwankungen nicht ausschließen.
  \item \textbf{Selbstauskunft:} [Wenn Befragungen eingesetzt wurden]
        Antworten basieren auf Selbsteinschaetzungen der Teilnehmenden.
\end{itemize}
